\chapter{Resultados}
%não se esqueça de definir uma label única para utilizar no comando \ref
\label{chap:resultados}

É aqui que você finalmente mostra o fruto de todo o seu esforço, suas noites mal dormidas e aquela planilha que travou o excel 7 vezes. Se tudo correu bem, você apresenta resultados lindos, organizados e que reforçam sua hipótese. Se não correu tão bem (geralmente é o caso)… você apresenta o que deu pra salvar com uma legenda bem convincente. 

O mais importante é escrever de forma segura, com honestidade intelectual, atenção e rigor científico, mesmo que por dentro você esteja pensando \enquote{que a banca tenha piedade.}

%Please add the following packages if necessary:
%\usepackage{booktabs, multirow} % for borders and merged ranges
%\usepackage{soul}% for underlines
%\usepackage{xcolor,colortbl} % for cell colors
%\usepackage{changepage,threeparttable} % for wide tables
%If the table is too wide, replace \begin{table}[!htp]...\end{table} with
%\begin{adjustwidth}{-2.5 cm}{-2.5 cm}\centering\begin{threeparttable}[!htb]...\end{threeparttable}\end{adjustwidth}
\begin{table}[!htp]
\centering
\caption{Exemplo de tabela.}
\label{tab:tabela1}
\Large
\scriptsize

{\normalsize
\begin{tabular}{|l|c|c|c|c|c|c|c|c|}
\hline
\textbf{Defesa} &
\rotatebox{90}{\textbf{FGSM}} & 
\rotatebox{90}{\textbf{DeepFool    }} & 
\rotatebox{90}{\textbf{JSMA}} & 
\rotatebox{90}{\textbf{C\&W}} & 
\rotatebox{90}{\textbf{BIM}} & 
\rotatebox{90}{\textbf{PGD}} & 
\rotatebox{90}{\textbf{L-BFGS}} & 
\rotatebox{90}{\textbf{UAP}} \cr
\hline
Treinamento Adversarial & X & X & X & & & & & \\
\hline
Compressão de Dados como Defesa & X & & & X & & & & O \\
\hline
MagNet & X & X & & O & X & & & \\
\hline
Destilação Defensiva & X & & X & O & & & & \\
\hline
Aleatorização & X & X & & X & & & & \\
\hline
Feature Squeezing & X & X & X & X & X & & & \\
\hline
Regularização / Máscara do Gradiente & X & & X & X & & X & X & \\
\hline
Redes de Parseval & X & & & & & & & \\
\hline
Deep Contractive Networks & & & & & & & X & \\
\hline
DeepCloak & X & & & & & & & \\
\hline
SafetyNet & X & X & & & X & & X & \\
\hline
Defesa Baseada em Foveação & X & & & & & & X & \\
\hline
PRN & & & & & & & & X \\
\hline
Sub-rede de Detecção & X & X & & & X & & & \\
\hline
Super-Resolução de Imagem & X & X & & X & X & & & \\
\hline
\end{tabular}
}

%Adicionar um trecho descrevendo a razão para não estarem na tabela todas as formas de defesa/ataque: não haver uma correlação possível.

\fonteautor
\end{table}


Ou até mesmo colocar na horizontal pra caber tudo que você precisa, apesar de ficar esquisito: quadro \ref{quad:quadro1}. Tabelas, como a tabela \ref{tab:tabela1} são usadas para exibir dados numéricos, enquanto quadros são utilizados para apresentar informações textuais, com ou sem a presença de dados numéricos.

\begin{landscape}
\begin{quadro}[!htp]
\centering
\caption{Exemplo de quadro.}
\label{quad:quadro1}
\scriptsize

\begin{tabular}{|p{2.5cm}|p{5cm}|p{8cm}|p{6cm}|}
\hline
\textbf{Recurso} & \textbf{URL} & \textbf{Ataques} & \textbf{Defesas} \\
\hline
AdvBox & \url{https://github.com/advboxes/AdvBox} & C\&W (L2 norm), DeepFool, FGSM, BIM, L-BFGS, JSMA, One Pixel Attack & Feature Squeezing, Gaussian Augmentation, Label Smoothing, Spatial Smoothing, Thermometer Encoding \\
\hline
Advertorch & \url{https://github.com/BorealisAI/advertorch} & C\&W (L2 norm), DeepFool, EAD (L1 norm), PGD, BIM, JSMA, L-BFGS, One Pixel Attack & Data Compression, Label Smoothing, Feature Squeezing \\
\hline
AISafety & \url{https://openi.pcl.ac.cn/OpenI/AISafety} & FGSM, AutoPGD, Boundary Attack, BIM, L-BFGS (BLB), C\&W (L2 norm), AOA, DeepFool, DIM, EAD, JSMA, PGD, SimBA, Square Attack, TIM, UAP, ZOO & Adversarial Training, Randomization \\
\hline
ART (Adversarial Robustness Toolbox) & \url{https://adversarial-robustness-toolbox.readthedocs.io/} & AutoPGD, Auto-CG, Boundary Attack, B\&B, C\&W, DeepFool, EAD, FGSM, JSMA, HSJA, BIM, PGD, NewtonFool, Shadow Attack, Wasserstein Attack, ZOO, Square Attack, UAP & Adversarial Training, Defensive Distillation \\
\hline
CleverHans & \url{https://github.com/cleverhans-lab/cleverhans#setting-up-cleverhans} & C\&W (L2 norm), FGSM, HSJA, PGD, BIM, DeepFool, EAD, L-BFGS, JSMA & - \\
\hline
DEEPSEC & \url{https://github.com/ryderling/DEEPSEC} & C\&W (L2 norm), FGSM, PGD, BIM, DeepFool, EAD, L-BFGS, JSMA, UAP & Defensive Distillation, Randomization, Gradient Regularization / Masking, Thermometer Encoding \\
\hline
Foolbox & \url{https://foolbox.readthedocs.io/en/stable/#} & FGSM, PGD, B\&B, Boundary Attack, DeepFool, ZOO & - \\
\hline
RobustBench & \url{https://robustbench.github.io/} & - & - \\
\hline
DI-2-FGSM Implementation & \url{https://github.com/cihangxie/DI-2-FGSM} & FGSM & - \\
\hline
NIPS 2017 Adversarial Defense Challenge & \url{https://github.com/cihangxie/NIPS2017_adv_challenge_defense} & - & Randomization \\
\hline
Facebook Archive Adversarial Defenses & \url{https://github.com/facebookarchive/adversarial_image_defenses} & FGSM, DeepFool, C\&W, UAP & Data Compression \\
\hline
Pixel Deflection Implementation & \url{https://github.com/iamaaditya/pixel-deflection} & - & Pixel Deflection \\
\hline
\end{tabular}

\fonteautor
\end{quadro}
\end{landscape}
